\documentclass[a4paper,twocolumn,11pt,unpublished]{quantumarticle}
\pdfoutput=1
\usepackage[utf8]{inputenc}
\usepackage[T1]{fontenc}
\usepackage[english]{babel}
\usepackage{amsmath,amssymb,amsthm,booktabs,graphicx,microtype}
\usepackage[numbers,sort&compress]{natbib}
\usepackage{hyperref,orcidlink,enumitem}
\hypersetup{allcolors=quantumviolet}
\graphicspath{{figures/}}
\newtheorem{theorem}{Theorem}
\newcommand{\argmax}{\operatorname*{arg\,max}}
\newcommand{\norm}[1]{\lVert#1\rVert}
\begin{document}
\title{Certified Training-Free Prescreening of Dissipators for Quantum Reservoir Computing}
\author[qar,mobile]{Markus Baumann}
\orcid{0009-0007-3575-1006}
\email{markus.baumann@campus.lmu.de}
\author[mobile]{Itamar Fink}
\author[mobile]{Johannes Wittmann}
\author[qar,mobile]{Claudia Linnhoff-Popien}
\author[qar,mobile]{Jonas Stein}
\affiliation[qar]{Quantum Applications and Research Laboratory, Ludwig-Maximilians-Universit\"at M\"unchen, Munich, Germany}
\affiliation[mobile]{Mobile and Distributed Systems Group, Ludwig-Maximilians-Universit\"at M\"unchen, Munich, Germany}
\begin{abstract}
For a fixed quantum reservoir, a practitioner normally chooses among a finite
hardware-compatible set of dissipators. Exhaustively generating trajectories,
training a readout, and testing every candidate is costly. We derive a
training-free score from first- and second-order derivatives of the driven
quantum channel and convert its approximation error into deterministic
selection guarantees. If all candidate scores have uniform error at most
$\delta$, the predicted winner has regret at most $2\delta$; a margin larger
than $2\delta$ certifies the exact winner; and interval elimination retains
every true optimum. An exact equal-decay-budget commuting quantum reservoir
shows that recent and old-memory tasks can prefer incompatible environments.
In an exhaustive three-qubit benchmark of 84 fully trained candidate
conditions grouped into 21 decisions, training only the predicted winner
recovers the empirical optimum in 18 decisions, saves 75\% of full candidate
trainings, and incurs mean regret $2.01\times10^{-4}$. A conservative
leave-one-processor-out interval shortlist retains the optimum in all 21
decisions while saving 15.5\% of full trainings. The claim is deliberately
bounded to a fixed architecture, task, and declared feasible family.
\end{abstract}
\maketitle

\section{Problem and score}
Open QRC uses fixed quantum dynamics and a trained linear readout. Dissipation
can create fading memory~\cite{Sannia2024}, while environmental organization
changes which temporal traces remain accessible~\cite{Baumann2026}. We ask the
bounded engineering question: given one architecture, one task, and a finite
feasible family $\mathcal C=\{C_1,\ldots,C_M\}$, which candidates deserve full
training?

Let
\begin{equation}
 r_t=\Phi_{u_t,C}r_{t-1},\qquad x_t=Rr_t,
\end{equation}
and at a reference input $u_0$ define
\begin{equation}
 A_C=\Phi_{u_0,C},\quad D_C=\partial_u\Phi_{u,C}|_{u_0},\quad A_Cr_{*,C}=r_{*,C}.
\end{equation}
For binary perturbations $u_t=u_0+\epsilon z_t$, channel differentiation gives
\begin{align}
 h_d(C)&=RA_C^dD_Cr_{*,C},\\
 q_{a,b}(C)&=RA_C^aD_CA_C^{b-a-1}D_Cr_{*,C},\qquad a<b.
\end{align}
These are the first-order delay and second-order cross-delay directions of the
quantum Volterra expansion~\cite{Hu2024}. Walsh orthogonality gives the order-two
feature covariance
\begin{equation}
 G_C=\epsilon^2HH^\top+\epsilon^4QQ^\top+\Sigma_{\rm shot}.
\end{equation}
For target $z_{t-a}z_{t-b}$ the training-free population score is
\begin{equation}
 \widehat S_{a,b}(C)=(\epsilon^2q_{a,b})^\top G_C^+(\epsilon^2q_{a,b}).
\end{equation}
The same construction gives delay scores from $h_d$. It incorporates the
chosen readout and shot covariance rather than equating a long-lived mode with
usable memory.

\begin{figure*}[t]
\centering
\includegraphics[width=.96\textwidth]{prescreen_pipeline.pdf}
\caption{The task enters a response-based score. Only candidates that survive
the score and error intervals require full task-specific training.}
\end{figure*}

\section{Finite-family certificate}
Let $S_i$ denote the true population score and $\widehat S_i$ its proxy.
\begin{theorem}[Regret, recovery, and safe elimination]
Assume $\max_i|S_i-\widehat S_i|\le\delta$. Let
$i^*\in\argmax_iS_i$ and $\widehat i\in\argmax_i\widehat S_i$. Then
\begin{equation}
 S_{i^*}-S_{\widehat i}\le2\delta.
\end{equation}
If the predicted winner margin exceeds $2\delta$, it is the unique true winner.
Moreover, discarding candidate $j$ whenever
\begin{equation}
 \widehat S_j+\delta<\max_i(\widehat S_i-\delta)
\end{equation}
retains every true maximizer. If the predicted top-$k$ boundary gap exceeds
$2\delta$, the complete predicted top-$k$ set is exact.
\end{theorem}
\begin{proof}
Because $\widehat i$ maximizes the proxy,
\begin{align}
 S_{i^*}-S_{\widehat i}
 &=(S_{i^*}-\widehat S_{i^*})
 +(\widehat S_{i^*}-\widehat S_{\widehat i})
 +(\widehat S_{\widehat i}-S_{\widehat i})\\
 &\le\delta+0+\delta.
\end{align}
For exact recovery and elimination, compare the lower interval endpoint of the
predicted winner with every competitor's upper endpoint. The identical
argument across the top-$k$ boundary proves the final statement.
\end{proof}
Thus training $k$ survivors instead of all $M$ candidates saves a fraction
$1-k/M$ of full task-training runs.

For the ridge-regularized score $s_\lambda(G,g)=g^\top(G+\lambda I)^{-1}g$,
positive semidefinite $G,\widehat G$, and errors
$\norm{G-\widehat G}_2\le\epsilon_G$ and
$\norm{g-\widehat g}_2\le\epsilon_g$, the resolvent identity gives
\begin{equation}
 |s_\lambda(G,g)-s_\lambda(\widehat G,\widehat g)|
 \le \frac{\epsilon_g(2\norm{\widehat g}_2+\epsilon_g)}{\lambda}
 +\frac{\norm{\widehat g}_2^2\epsilon_G}{\lambda^2}.
\end{equation}
This supplies $\delta$ from feature, covariance, shot, and truncation bounds.
If $\norm{A}\le\eta<1$, first- and second-order lag tails beyond $L$ obey
\begin{equation}
 r_{L,2}\le |\epsilon|B_1\frac{\eta^{L+1}}{1-\eta}
 +|\epsilon|^2B_2\frac{(L+1)\eta^L-L\eta^{L+1}}{(1-\eta)^2}
 +r_{\ge3}.
\end{equation}

\section{Why task dependence is unavoidable}
Consider two independently damped computational-basis modes with
$\Lambda=\mathrm{diag}(e^{-\gamma_1},e^{-\gamma_2})$, full readout, and a
nonzero affine input write. This is a commuting CPTP quantum reservoir. For
independent unit-variance inputs, exact delay capacity is
\begin{equation}
 \mathcal C_d=r_d^\top G^{-1}r_d,\quad
 (r_d)_i=b_i\lambda_i^d,\quad
 G_{ij}=\frac{b_ib_j}{1-\lambda_i\lambda_j}.
\end{equation}
Hold the total decay budget fixed at $\gamma_1+\gamma_2=1$. The balanced
allocation $(0.4,0.6)$ has capacity $0.265180$ at delay 2, versus $0.078103$
for $(0.05,0.95)$. At delay 10 the ordering reverses: $0.003620$ versus
$0.043885$. Hence no one of these equal-budget environments is optimal for both
tasks.

\begin{figure}[t]
\centering
\includegraphics[width=\columnwidth]{exact_counterexample.pdf}
\caption{Exact equal-budget construction. Balanced decay favors recent memory;
heterogeneous decay creates one slow mode and favors old memory.}
\end{figure}

\section{Exhaustive practitioner benchmark}
The decision benchmark uses three independently drawn three-qubit XX+Z
processors, seven delayed-product tasks, and four matched-gap dissipator
candidates, giving 21 architecture--task decisions and 84 candidate conditions.
Every candidate receives a full switched Lindblad simulation, readout fit, and
held-out test score. The proxy score is computed without task-specific fitting.

At candidate level, predicted and trained capacities correlate at $0.99711$
and have MAE $0.00451$. Table~\ref{tab:shortlist} reports the actual decision
trade-off. Training only the predicted winner recovers 18 of 21 empirical
winners and saves 75\% of full candidate trainings. Mean regret is
$0.0002006$, maximum regret $0.003204$, and the Wilson 95\% interval for winner
recovery is $[65.4\%,95.0\%]$. The selected candidate improves mean empirical
capacity by $0.02076$ over always-local selection, $0.001786$ over always using
the most collective candidate, and $0.01445$ over a random candidate in
expectation.

\begin{table}[t]
\centering
\caption{Shortlist performance over 21 decisions. Maximum regret is 0.003204,
0.003204, 0.001776, and 0 for top-1 through top-4.}
\label{tab:shortlist}
\begin{tabular}{lrrr}
\toprule
Shortlist & Winner hit & Mean regret & Saved\\
\midrule
Top 1 & 85.7\% & 0.0002006 & 75\%\\
Top 2 & 85.7\% & 0.0002006 & 50\%\\
Top 3 & 95.2\% & 0.0000846 & 25\%\\
All 4 & 100\% & 0 & 0\%\\
\bottomrule
\end{tabular}
\end{table}

\begin{figure*}[t]
\centering
\includegraphics[width=.95\textwidth]{prescreen_results_combined.pdf}
\caption{Left: response score versus fully trained held-out capacity. Right:
full candidate trainings saved versus empirical shortlist regret.}
\end{figure*}

A conservative leave-one-processor-out radius, set to the maximum absolute
residual on the other processors, retains the empirical winner in all 21 held-
out decisions, gives zero observed shortlist regret, keeps 3.38 of four
candidates on average, and saves 15.5\% of full trainings. This is an empirical
calibration procedure; the deterministic theorem is guaranteed only when the
supplied radius truly bounds the deployment family.

Supporting checks include 105 exact nonlinear conditions across finite $N=3$
and $N=4$ systems (correlation $0.9891$), 72 finite-shot conditions
(correlation $0.9954$), 20 mixed finite-difference kernel audits (maximum
relative error $1.26\times10^{-7}$), and a 16-condition phase-orientation study
(correlation $0.9989$).

\section{Practitioner protocol and boundary}
A practitioner should: freeze the QRC architecture and measurement budget;
declare a finite, completely positive, equal-budget, convergence-screened
family; represent the target in an orthogonal temporal basis; calculate the
necessary channel derivatives and shot-aware score; attach an analytic or
conservatively calibrated error radius; eliminate candidates by score
intervals; and fully train only the survivors.

The result does not claim the globally optimal Lindbladian, asymptotic scaling,
or generic quantum advantage over classical reservoirs. A midpoint Liouvillian
gap is a feasibility diagnostic, not a proof of uniform switched-input
contractivity. The contribution is the certified bounded decision rule that a
QRC practitioner actually needs.

\section{Conclusion}
Within a fixed architecture and declared feasible dissipator family,
differentiated-channel scores provide training-free task-dependent
prescreening. Uniform score error yields deterministic regret and recovery
guarantees, while exhaustive tests show that most full candidate trainings can
be removed with negligible observed regret. Environmental coupling is thereby
a programmable design layer, but the scientific claim remains explicit and
bounded.

\begin{thebibliography}{9}
\bibitem{Dambre2012} J.~Dambre, D.~Verstraeten, B.~Schrauwen, and S.~Massar,
``Information processing capacity of dynamical systems,'' \emph{Scientific
Reports} \textbf{2}, 514 (2012).
\bibitem{Sannia2024} A.~Sannia, R.~Mart\'inez-Pe\~na, M.~C.~Soriano,
G.~L.~Giorgi, and R.~Zambrini, ``Dissipation as a resource for quantum
reservoir computing,'' \emph{Quantum} \textbf{8}, 1291 (2024).
\bibitem{Hu2024} F.~Hu \emph{et al.}, ``Overcoming the coherence time barrier in
quantum machine learning on temporal data,'' \emph{Nature Communications}
\textbf{15}, 7491 (2024).
\bibitem{Baumann2026} M.~Baumann \emph{et al.}, ``The organization of
environmental coupling shapes what quantum reservoirs remember,''
arXiv:2608.14181 (2026).
\end{thebibliography}
\end{document}
